\documentclass[11pt]{article}

\usepackage[margin=1in]{geometry}
\usepackage[T1]{fontenc}
\usepackage[utf8]{inputenc}
\usepackage{lmodern}
\usepackage{microtype}
\usepackage{amsmath,amssymb,amsthm}
\usepackage[numbers,sort&compress]{natbib}
\usepackage[colorlinks=true,linkcolor=blue!60!black,citecolor=blue!60!black,urlcolor=blue!60!black]{hyperref}

\title{\textbf{Canonical Quaternion Lifting for Multichannel RF and Sensing Signals}}
\author{John G. Van Geem\\RQM Technologies}
\date{Version 0.1.0 \ Draft dated 2026-04-23}

\begin{document}
\maketitle

\begin{abstract}
This draft paper defines a canonical quaternion lifting scaffold for multichannel RF and sensing signals, with an emphasis on publication-safe package structure rather than finalized technical claims. It outlines how aligned channel tuples may be lifted into quaternion-valued representations, specifies placeholder conventions for orientation and sign choices, and records interoperability surfaces needed for reproducible downstream indexing. The package is intentionally layout-only at this stage and does not report benchmark results, empirical performance figures, or publication metadata beyond draft status.
\end{abstract}

\noindent\textbf{Authoring note:} This file is the source of truth for this draft paper package.

\section{Draft scaffold notice}
This manuscript is intentionally in scaffold form only.

\section{Planned content}
Full technical sections will be authored in a future revision.

\bibliographystyle{plainnat}
\bibliography{references}

\end{document}
